% !TeX encoding = UTF-8 Unicode
% !TeX program = xelatex
% !TeX spellcheck = sk_SK
%
\documentclass[11pt,a4paper%,twoside
]{book}
\usepackage{xltxtra}
\usepackage{amssymb}
\usepackage{amsmath}
\usepackage{fontspec}
\usepackage{unicode-math}
\newcommand{\noteDynamics}[2][1.3]{\raisebox{0.05ex}{\fontspec[Scale=#1,Path=fonts/]{emmentaler-11.otf}{#2}}}
\newcommand{\noteCommonTime}{\makebox[2ex]{\raisebox{0.8ex}{\fontspec[Scale=1.3,Path=fonts/]{emmentaler-11.otf}{}}}}
\newcommand{\noteCutTime}{\makebox[2ex]{\raisebox{0.8ex}{\fontspec[Scale=1.3,Path=fonts/]{emmentaler-11.otf}{}}}}
\newcommand{\noteSharp}{\makebox[0.8ex]{\raisebox{0.75ex}{\fontspec[Scale=1,Path=fonts/]{emmentaler-11.otf}{}}}}
\newcommand{\noteFlat}{\makebox[0.8ex]{\raisebox{0.50ex}{\fontspec[Scale=1,Path=fonts/]{emmentaler-11.otf}{}}}}
\setmainfont[Path=fonts/,
Numbers={Lining},
Ligatures={Common,TeX},
BoldFont=LinLibertineOB.otf,
ItalicFont=LinLibertineOI.otf,
BoldItalicFont=LinLibertineOBI.otf
]{LinLibertineO.otf}
\setmathfont[Path=fonts/]{texgyreschola-math.otf}
\newfontfamily{\noligatures}[Ligatures=NoCommon,Path=fonts/]{LinLibertineO.otf}
\usepackage{booktabs}
\usepackage{rotating}
\usepackage{polyglossia}
\setdefaultlanguage{slovak}
\setotherlanguage{english}
\usepackage{xcolor}
\usepackage{natbib}
\usepackage{hyperref}
\definecolor{printblue}{rgb}{0,0,0.8}
\definecolor{proofreadcolor}{rgb}{0.0157,0.5176,0.2666}
\definecolor{attentioncolor}{rgb}{0.8235,0.2666,0}
\hypersetup{bookmarksnumbered=true,
bookmarksopen=false,
pdfpagemode=UseNone,
hypertexnames=false,
colorlinks=true,
urlcolor=printblue,
filecolor=printblue,
linkcolor=printblue,
citecolor=printblue,
anchorcolor=printblue,
hypertexnames=false,  
pdftitle={Základy cvičenia na klavíri},
pdfauthor={Chuan C. Chang},
pdfsubject={Metódy a techniky cvičenia na klavíri},
pdfkeywords={klavír, piano, cvičenie, technika, ladenie klavíra},
unicode=true,
pdfstartview=FitH}
\usepackage[left=1in,top=1in,right=1in,bottom=1in%,bindingoffset=0.5in
]{geometry}
%\usepackage{setspace}
%\setstretch{1.15}
\usepackage{multicol}
\usepackage{enumerate}
\usepackage{makeidx} 
\makeindex
\setcounter{tocdepth}{5}
\setcounter{secnumdepth}{4}
\renewcommand\thechapter{}
\renewcommand\thesection{\Roman{section}}
\renewcommand\thesubsection{\thesection.\arabic{subsection}}
\renewcommand\thesubsubsection{\thesubsection.\alph{subsubsection}}
\usepackage{array}
\newcolumntype{C}[1]{>{\centering\let\newline\\\arraybackslash\hspace{0pt}}m{#1}}
\newcolumntype{L}[1]{>{\raggedright\let\newline\\\arraybackslash\hspace{0pt}}m{#1}}
\newcolumntype{R}[1]{>{\raggedleft\let\newline\\\arraybackslash\hspace{0pt}}m{#1}}
\newcommand{\TODO}[1]{{\color{red} #1}}
\newcommand{\Proofread}[1]{{\color{proofreadcolor} #1}}
\newcommand{\NeedRevision}[1]{{\color{attentioncolor} #1}}
\usepackage{fancyhdr}
\begin{document}

\textbf{III. VYBRANÉ TÉMY Z KLAVÍRNEJ PRAXE}

\textbf{1. Tón, Rytmus, Legato, Staccato}

\begin{enumerate}
\def\labelenumi{\alph{enumi}.}
\item
  \textbf{Čo je to „dobrý tón``?}
\end{enumerate}

\textbf{Základné stlačenie kláves.} Základné stlačenie kláves sa musí naučiť každý klavirista. Bez neho nič iné nebude mať zmysluplný význam -- Tádž Mahal z hlinených tehál a slamy nepostavíte. \emph{\textbf{Stlačenie kláves sa skladá z 3 hlavných komponentov: stlačenia klávesu smerom nadol, podržania a zdvihnutia klávesu.}} Môže sa to zdať ako triviálne jednoduchá vec na naučenie, ale nie je to tak a väčšina učiteľov klavíra má problém naučiť svojich študentov správne stlačenie klávesu. Ťažkosti vznikajú najmä preto, že mechanika stlačenia klávesu nebola nikde dostatočne vysvetlená; preto budú tieto vysvetlenia hlavnými témami týchto odsekov.

Úder smerom nadol je to, čo spočiatku vytvára zvuk klavíra; pri správnom pohybe musí byť čo najrýchlejší, ale s kontrolou hlasitosti. Toto ovládanie nie je jednoduché, pretože v časti o gravitačnom poklese sme zistili, že rýchlejší úder smerom nadol vo všeobecnosti znamená hlasnejší zvuk. Rýchlosť dáva note presné načasovanie; bez tejto rýchlosti sa načasovanie začiatku noty stáva nedbalou záležitosťou. Preto, či už je hudba pomalá alebo rýchla, úder smerom nadol musí byť v podstate rýchly. \emph{\textbf{Tieto požiadavky na rýchly úder, ovládanie hlasitosti a mnohé ďalšie, s ktorými sa čoskoro stretneme, nás privádzajú k najdôležitejšiemu princípu učenia sa hry na klavíri -- citlivosti prstov. Prst musí byť schopný vnímať a vykonávať mnoho požiadaviek, kým zvládnete základné údery klávesov.}} Aby ste mohli ovládať hlasitosť, úder smerom nadol by mal pozostávať z 2 častí; počiatočná silná zložka na prerušenie trenia/zotrvačnosti klávesu a spustenie jeho pohybu a druhý komponent s vhodnou silou pre požadovanú hlasitosť. Návrh „hrať hlboko do klávesov`` je dobrý v tom zmysle, že úder nadol sa nesmie spomaliť; musí sa zrýchliť, keď sa dostanete na spodok, aby ste nikdy nestratili kontrolu nad kladivkom.

\emph{\textbf{Tento dvojčasťový pohyb je obzvlášť dôležitý pri hre pianissimo.}} V dobre regulovanom koncertnom krídle je trenie takmer nulové a zotrvačnosť systému je nízka. Vo všetkých ostatných klavíroch (ktoré tvoria 99 \% všetkých klavírov) existuje trenie, ktoré treba prekonať, najmä pri prvom spustení úderu nadol (trenie je najvyššie, keď je pohyb nulový), a v systéme existuje množstvo nerovnováh, ktoré vytvárajú zotrvačnosť. Za predpokladu, že klavír má správny zvuk, môžete hrať veľmi jemné pianissimo tak, že najprv prerušíte trenie/zotrvačnosť a potom vykonáte jemný úder. Tieto dve zložky sa musia plynule spojiť, aby to pre pozorovateľa vyzeralo ako jeden pohyb, pričom mäso prstov funguje ako tlmiče nárazov. Požadovaný rýchly úder nadol znamená, že sval prsta musí mať vysoký podiel rýchlych svalov (pozri časť 7.a nižšie). To sa dosahuje dlhším časovým cvičením rýchleho pohybu (približne rok) a vyhýbaním sa silovým cvičeniam; preto je tvrdenie, že technika hry na klavír vyžaduje silu prstov, absolútne nesprávne. Musíme si pestovať rýchlosť a citlivosť prstov.

Podržanie klávesu je nevyhnutné na udržanie kladivkového úderu pomocou spätného rázu a na presné ovládanie trvania tónu. Bez podržania sa kladivkové úderové údery môžu kývať, produkovať cudzie zvuky, spôsobovať problémy s opakovanými tónmi, trilkovaním atď. Začiatočníci budú mať problém s plynulým prechodom medzi úderom nadol a podržaním. \emph{\textbf{Počas podržania netlačte na kláves v snahe „zatlačiť hlboko do klavíra``; gravitácia je dostatočná na to, aby kláves zostal stlačený.}} Dĺžka podržania je to, čo riadi farbu a výraz; preto je dôležitou súčasťou hudby.

Zdvih spôsobí, že tlmič dopadne na struny a ukončí zvuk. Spolu s podržaním určuje trvanie noty. Rovnako ako pri údere nadol, aj zdvih musí byť rýchly, aby sa trvanie noty presne ovládalo. Preto sa musí klavirista vedome snažiť rozvíjať rýchle svaly v extenzorových svaloch, rovnako ako sme to urobili s flexorovými svalmi pri údere nadol. Najmä pri rýchlej hre mnohí študenti na zdvih úplne zabudnú, čo vedie k nedbalej hre. Séria môže skončiť pozostávajúca zo staccata, legata a prekrývajúcich sa tónov. Rýchle paralelné série môžu nakoniec znieť, akoby sa hrali pedálom.

\emph{\textbf{Presným ovládaním všetkých 3 komponentov základného úderu kláves si udržiavate úplnú kontrolu nad klavírom, konkrétne nad kladivkom a tlmičom, a toto ovládanie je potrebné pre autoritatívnu hru.}} Tieto komponenty určujú povahu každej noty. Teraz vidíte, prečo je rýchly úder smerom nadol a rovnako rýchly zdvih taký dôležitý, najmä počas pomalej hry. Pri normálnej hre sa zdvih predchádzajúcej noty zhoduje so zdvihom nasledujúcej noty smerom nadol. Pri staccato a legato (časť c) a rýchlej hre (7.i) musíme všetky tieto komponenty upraviť a budú prediskutované samostatne. Ak ste tieto komponenty nikdy predtým necvičili, začnite cvičiť so všetkými 5 prstami, C až G, ako to robíte pri hraní stupnice, a aplikujte komponenty na každý prst, HS. Ak chcete precvičiť extenzorové svaly, môžete rýchly zdvih prehnať. Snažte sa všetky nehrajúce prsty udržať na klávesoch zľahka. Keď sa snažíte zrýchliť údery smerom nadol a nahor, hrajúc približne jednu notu za sekundu, môžete začať vytvárať stres; v takom prípade cvičte, kým stres nedokážete eliminovať. Najdôležitejšia vec, ktorú si treba zapamätať o výdrži, je, že po rýchlom stlačení nadol sa musíte počas výdrže okamžite uvoľniť. Inými slovami, \emph{\textbf{musíte si precvičiť oboje, rýchlosť úderu aj rýchlosť relaxácie}}. Potom postupne zrýchľujte hru; v tejto chvíli nie je potrebné hrať rýchlo. Stačí dosiahnuť pohodlné tempo. Teraz urobte to isté s akoukoľvek pomalou hudbou, ktorú viete hrať, napríklad s 1. vetou Beethovenovej Sonáty mesačného svitu, HS. Ak ste to nikdy predtým nerobili, bude hranie HT spočiatku veľmi neobratné, pretože teraz musíte koordinovať toľko komponentov v oboch rukách. S cvičením sa však hudba bude zlepšovať, získate oveľa väčšiu kontrolu nad výrazom a mali by ste nadobudnúť pocit, že teraz môžete hrať muzikálnejšie. Nemali by ste mať žiadne vynechané alebo nesprávne noty, všetky noty by mali byť rovnomernejšie a všetky výrazové značky môžete hrať s väčšou presnosťou. Výkony budú zo dňa na deň konzistentné a technika sa bude zlepšovať rýchlejšie. Bez dobrého základného úderu kláves sa môžete dostať do problémov, keď hráte na rôznych klavíroch alebo na klavíroch, ktoré nie sú v dobrej regulácii, a hudba sa môže často zhoršiť, keď budete cvičiť viac, pretože si môžete osvojiť zlé návyky, ako je nepresné načasovanie. Samozrejme, celý proces opísaný v tomto jednom odseku môže trvať týždne alebo dokonca mesiace.

\textbf{Tón: Jedna verzus viacero nôt, Pianissimo, Fortissimo. \emph{Tón je kvalita zvuku; či je súčet všetkých vlastností zvuku vhodný pre danú hudbu}}. Existuje kontroverzia o tom, či klavirista dokáže ovládať „tón`` jednej noty na klavíri. Ak by ste si sadli za klavír a hrali jednu notu, zdá sa byť takmer nemožné zmeniť tón okrem vecí ako staccato, legato, hlasno, ticho atď. Na druhej strane niet pochýb o tom, že rôzni klaviristi produkujú rôzne tóny. Dvaja klaviristi môžu hrať tú istú skladbu na tom istom klavíri a produkovať hudbu s veľmi odlišnou tonálnou kvalitou. Väčšinu tohto zdanlivého rozporu možno vyriešiť starostlivým definovaním toho, čo znamená „tón``. Napríklad veľkú časť tónových rozdielov medzi klaviristami možno pripísať konkrétnym klavírom, ktoré používajú, a spôsobu, akým sú tieto klavíry regulované alebo ladené. Ovládanie tónu jednej noty je pravdepodobne len jedným aspektom mnohostranného a zložitého problému. Preto najdôležitejším rozdielom, ktorý musíme na začiatku urobiť, je, či hovoríme o jednej note alebo o skupine tónov. Keď počujeme rôzne tóny, väčšinou počúvame skupinu nôt. V takom prípade sa rozdiely v tónoch ľahšie vysvetľujú. Tón vzniká väčšinou ovládaním nôt vo vzťahu k sebe navzájom. To takmer vždy závisí od presnosti, kontroly a hudobného obsahu. \emph{\textbf{Preto je tón hlavne vlastnosťou skupiny nôt a závisí od hudobnej citlivosti klaviristu.}}

\emph{\textbf{Je však tiež zrejmé, že tón jednej noty môžeme ovládať niekoľkými spôsobmi.}} Môžeme ho ovládať pomocou jemných a tlmiacich pedálov. Harmonický obsah (počet alikvotných tónov) môžeme tiež zmeniť hraním hlasnejšie alebo tichšie. Jemný pedál mení tón alebo zafarbenie tónu znížením okamžitého zvuku v porovnaní s dozvukom. Keď sa na strunu udiera s väčšou silou, generuje sa viac harmonických. Keď teda hráme jemne, máme tendenciu produkovať zvuk obsahujúci silnejšie základné tóny. Pod určitou hlasitosťou však môže byť nedostatočná energia na vybudenie základného tónu a môžete vybudiť len niektoré postupné vlny s vyššou frekvenciou, trochu podobné flautandu na husliach (zotrvačnosť strún pôsobí ako prst vo flautande). Preto niekde medzi PP a FF existuje optimálna sila úderu, ktorá maximalizuje základný tón. Tlmiaci pedál tiež mení zafarbenie tónu tým, že umožňuje vibrácie na nehraných strunách.

\emph{\textbf{Tón alebo zafarbenie tónu môže ladič ovládať nastavením kladivka alebo odlišným ladením.}} Tvrdšie kladivko vytvára brilantnejší tón (väčší harmonický obsah) a kladivko s plochou úderovou plochou vytvára drsnejší tón (viac vysokofrekvenčných harmonických). Ladič môže meniť natiahnutie alebo ovládať mieru rozladenia medzi unisonami. Do určitého bodu má väčšie natiahnutie tendenciu vytvárať jasnejšiu hudbu a nedostatočné natiahnutie môže viesť k nevýraznému zvuku klavíra. Pri rozladení v rozsahu sympatických vibrácií budú všetky struny noty dokonale naladené (vibrovať na rovnakej frekvencii), ale budú na seba reagovať odlišne. Napríklad nota môže „spievať``, čo môže byť dozvuk, ktorého hlasitosť kolíše. Žiadne dve struny nie sú nikdy identické, takže možnosť identického ladenia jednoducho neexistuje.

Nakoniec sa dostávame k zložitej otázke: \emph{\textbf{môžete zmeniť tón jednej noty ovládaním úderu nadol?}} Väčšina argumentov o ovládaní tónu sa sústreďuje na vlastnosť voľného letu kladivka predtým, ako udrie na struny. Odporcovia (ovládania tónu jednotlivými notami) tvrdia, že keďže kladivko je vo voľnom lete, dôležitá je iba jeho rýchlosť, a preto tón nie je ovládateľný pre notu hranú s určenou hlasitosťou. Predpoklad voľného letu však nikdy nebol dokázaný, ako teraz uvidíme. \emph{\textbf{Jeden faktor ovplyvňujúci tón je ohybnosť drieku kladivka.}} Pri hlasnom tóne sa driek môže výrazne ohnúť, keď sa kladivko vypustí do voľného letu. V takom prípade môže mať kladivko väčšiu efektívnu hmotnosť ako svoju pôvodnú hmotnosť, keď udrie do strún. Je to preto, že sila F kladiva na struny je daná vzťahom F = M \textbf{a}, kde M je hmotnosť kladiva a \textbf{a} je jeho spomalenie pri náraze na struny. Kladná ohybnosť pridáva ďalšiu silu, pretože keď sa ohybnosť po uvoľnení zdviháka vráti do pôvodného stavu, tlačí kladivko dopredu; keď sa F zvýši, nezáleží na tom, či sa zvýši M alebo \textbf{a}, účinok je rovnaký. Meranie \textbf{a} je však ťažšie ako M (napríklad väčšie M môžete ľahko simulovať použitím ťažšieho kladivka), takže v tomto prípade zvyčajne hovoríme, že „efektívna hmotnosť`` sa zvýšila, aby sme si ľahšie predstavili vplyv väčšieho F na reakciu strún. V skutočnosti však kladná ohybnosť zvyšuje \textbf{a}. Pri note hranej staccato môže byť ohybnosť negatívna v čase, keď kladivko udrie na struny, takže rozdiel tónov medzi „hlbokou`` hrou a staccatom môže byť značný. Tieto zmeny v efektívnej hmotnosti určite zmenia rozloženie alikvotných tónov a ovplyvnia tón, ktorý počujeme. \emph{\textbf{Keďže driek nie je 100 \% pevný, vieme, že vždy existuje konečná flexibilita. Jedinou otázkou je, či je dostatočná na ovplyvnenie tónu, ako ho počujeme.}} Je to takmer určite preto, že driek kladivka je relatívne ohybný kus dreva. Ak je to pravda, potom by mal byť tón nižších tónov pri ťažších kladivkách lepšie ovládateľný, pretože ťažšie kladivká spôsobia väčšie ohyb. Hoci by sa dalo očakávať, že ohyb bude zanedbateľný, pretože kladivko je také ľahké, kĺb je veľmi blízko puzdra príruby kladivka, čo vytvára obrovský pákový efekt. Argument, že kladivko je príliš ľahké na to, aby vyvolalo ohyb, nie je platný, pretože kladivko je dostatočne masívne na to, aby udržalo všetku kinetickú energiu potrebnú na vydanie aj tých najhlasnejších zvukov. To je veľa energie!

Všimnite si, že vzdialenosť uvoľnenia kladivka je len niekoľko milimetrov a táto vzdialenosť je mimoriadne dôležitá pre tón. Takáto malá vzdialenosť uvoľnenia naznačuje, že kladivko je navrhnuté tak, aby pri náraze na strunu dosahovalo zrýchlenie. Kladivko nie je vo voľnom lete po uvoľnení zdviháka, pretože počas prvých niekoľkých milimetrov po uvoľnení je kladivko zrýchľované obnovením ohybu drieku. Vzdialenosť uvoľnenia je najmenšia ovládateľná vzdialenosť, ktorá dokáže udržať zrýchlenie bez akejkoľvek možnosti zablokovania kladivka na strunách, pretože zdvihák by sa nemohol uvoľniť. Táto vzdialenosť uvoľnenia vysvetľuje štyri inak záhadné fakty: (i) obrovskú energiu, ktorú také ľahké kladivko dokáže preniesť na struny, (ii) zníženie kvality tónu (alebo ovládania), keď je vzdialenosť uvoľnenia príliš veľká, (iii) kritickú závislosť zvukového výstupu a ovládania tónu od hmotnosti a veľkosti kladivka a (iv) klikavý zvuk, ktorý klavír vydáva, keď sa opotrebuje puzdro drieku kladivka (klasickým príkladom je klikajúce teflónové puzdro). Klikanie je zvuk, keď sa puzdro zacvakne späť, keď sa zdvihák uvoľní a prevezme pôsobenie ohybnosť drieku -- bez uvoľnenia ohybnosti nie je potrebná žiadna sila na zacvaknutie puzdra späť, preto bez ohybnosti nebude počuť žiadne kliknutie. Keďže kliknutie je počuť aj pri mierne tichých zvukoch, driek je ohnutý pri všetkých zvukoch okrem tých najtichších.

Tento scenár má dôležité dôsledky aj pre klaviristu (nielen pre ladiča klavíra). Znamená to, že tón jednej noty je možné ovládať. Tiež nám to hovorí, ako ho ovládať. Po prvé, pri zvukoch PPP je zanedbateľné ohýbanie a máme do činenia s iným tónom ako pri hlasnejších zvukoch. Klaviristi vedia, že na hranie PP musíte stláčať nadol konštantnou rýchlosťou - všimnite si, že toto minimalizuje ohyb, pretože pri uvoľnení nedochádza k žiadnemu zrýchleniu. Pri hraní pianissima chcete minimalizovať ohyb, aby ste minimalizovali efektívnu hmotnosť kladivka. Po druhé, pre maximálny ohyb by sa mal zdvih zrýchliť v spodnej časti. Toto dáva veľký zmysel: „hlboký tón`` sa vytvára opretím sa o klavír a silným stlačením, a to aj pri tichých zvukoch. Presne takto maximalizujete ohyb, čo je ekvivalentné použitiu väčšieho kladivka. Táto informácia je tiež kľúčová pre klavírneho technika. Znamená to, že optimálna veľkosť kladivka je taká, ktorá je dostatočne malá, aby flexia bola nulová niekde okolo PP, ale dostatočne veľká, aby flexia bola významná od mf. Ide o veľmi šikovné mechanické usporiadanie, ktoré umožňuje použitie relatívne malých kladiviek, ktoré umožňujú rýchle opakovania a stále dokážu prenášať maximálne množstvo energie na struny. To znamená, že je chybou prejsť na väčšie kladivká na dosiahnutie väčšieho zvuku, pretože stratíte rýchlosť opakovania a kontrolu tónu. Existencia flexibility drieku kladivka je teraz dobre známa („Päť prednášok o akustike klavíra``).

Dá sa na klavíri počuť rozdiel v tóne jednej noty hraním iba jednej noty? Zvyčajne nie; väčšina ľudí nie je dostatočne citlivá na to, aby tento rozdiel počula pri väčšine klavírov. Budete potrebovať klavír Steinway B alebo lepší a tento rozdiel môžete začať počuť (ak to otestujete s niekoľkými klavírmi postupne vyššej kvality) už pri nižších tónoch. Avšak pri hraní skutočnej hudby je ľudské ucho úžasne citlivé na to, ako kladivko dopadá na struny, a rozdiel v tóne sa dá ľahko počuť. Je to podobné ako pri ladení: väčšina ľudí (vrátane väčšiny klaviristov) len ťažko počuje rozdiel medzi superladením a bežným ladením hraním jednotlivých tónov alebo dokonca testovaním intervalov. Prakticky každý klavirista však dokáže počuť rozdiel v kvalite ladenia hraním skladby svojej obľúbenej hudby. Môžete si to sami predviesť. Zahrajte jednoduchú skladbu dvakrát rovnakým spôsobom, okrem úderu. Najprv hrajte s váhou paže a „hlboko zatlačte`` do klavíra a uistite sa, že pokles kláves sa zrýchľuje úplne nadol (správny základný úder). Potom to porovnajte s hudbou, keď stlačíte plytko, takže dôjde k úplnému poklesu klávesu, ale v spodnej časti nie je žiadne zrýchlenie. Možno budete musieť trochu precvičiť, aby ste sa uistili, že prvýkrát nebudete hlasnejší ako druhýkrát. Pri druhom režime hry by ste mali počuť horšiu kvalitu tónu. V rukách skvelých klaviristov môže byť tento rozdiel dosť veľký. Samozrejme, vyššie sme hovorili o tom, že tón je najsilnejšie ovládaný tým, ako hráte po sebe idúce tóny, takže hranie hudby na testovanie účinku jednotlivých tónov zjavne nie je najlepší spôsob. Je to však najcitlivejší test.

\textbf{Pianissimo}: Videli sme, že pre PPP potrebujete presný základný úder a rýchle uvoľnenie. Dôležitý je pocit z klávesov vankúšikmi prstov. Vo všeobecnosti by ste mali vždy cvičiť s jemným úderom, kým nezvládnete pasáž, potom pridať mf alebo FF alebo čokoľvek, čo je potrebné, pretože hranie s jemným úderom je najťažšie rozvíjať. Nedochádza k zrýchleniu úderu nadol ani k ohýbaniu drieku kladivka, ale spätný úder musí byť kontrolovaný (kláves stlačený a podržaný). \emph{\textbf{Najdôležitejšími faktormi pre PPP sú správna regulácia (najmä minimálne uvoľnenie, zvuk kladivka a správna hmotnosť kladivka). Snažiť sa kultivovať techniku PPP bez riadnej údržby klavíra je márne.}} V núdzi (počas hrania s neuspokojivým klavírom) môžete skúsiť jemný pedál so vzpriameným alebo čiastočne jemným pedálom na krídle. PPP je na väčšine digitálnych nástrojov náročné, pretože chod kláves je horší a po približne 5 rokoch používania sa rýchlo zhorší. Ale akustický klavír, ktorý nebol udržiavaný, môže byť oveľa horší.

\textbf{Fortissimo} je o prenose váhy na klavír. To znamená, že sa telo nakloní dopredu, aby bolo ťažisko bližšie ku klávesnici, a \emph{\textbf{hrá sa z ramien}}. Pri hraní na FF nepoužívajte iba ruky alebo paže. Opäť platí, že relaxácia je dôležitá, aby ste neplytvali energiou, aby ste umožnili maximálnu rýchlosť úderu nadol a aby správna sila mohla byť nasmerovaná iba tam, kde je potrebná. \emph{\textbf{Ak chcete, aby sa pasáž hrala na FF, cvičte bez FF, kým si ju nezvládnete, a potom pridajte FF.}}

Stručne povedané, tón je primárne výsledkom jednotnosti a kontroly hry a závisí od hudobnej citlivosti hráča. \emph{\textbf{Ovládanie tónu je zložitá záležitosť zahŕňajúca}} \emph{\textbf{každý faktor, ktorý mení povahu zvuku, a videli sme, že existuje mnoho spôsobov, ako zmeniť zvuk klavíra.}} Všetko to začína tým, ako je klavír regulovaný. Každý klavirista môže ovládať tón mnohými spôsobmi, napríklad hraním hlasno alebo ticho, alebo zmenou rýchlosti. Napríklad hraním hlasnejšie a rýchlejšie môžeme vytvoriť hudbu pozostávajúcu prevažne z okamžitého zvuku; naopak, pomalšia a tichšia hra vytvorí tlmený efekt s použitím väčšieho množstva dozvuku. A existuje nespočetné množstvo spôsobov, ako začleniť pedál do vašej hry. Videli sme, že tón jednej noty sa dá ovládať, pretože driek kladivka je ohybný. Veľký počet premenných zaisťuje, že každý klavirista bude produkovať iný tón.

\begin{enumerate}
\def\labelenumi{\alph{enumi}.}
\setcounter{enumi}{1}
\item
  \textbf{Čo je rytmus?} (Beethovenova Búrka, Op. 31, \#2, Appassionata, Op. 57) 
\end{enumerate}
\emph{\textbf{Rytmus je (opakujúci sa) časový rámec hudby.}} Keď čítate o rytme (pozri Whiteside), často sa javí ako tajomný aspekt hudby, ktorý dokáže vyjadriť iba „vrodený talent``. Alebo ho možno musíte cvičiť celý život, ako bubeníci. Najčastejšie \emph{\textbf{je však správny rytmus jednoducho otázkou presného počítania, správneho čítania hudby, najmä taktových označení. Nie je to také jednoduché, ako sa zdá; ťažkosti často vznikajú, pretože väčšina označení rytmu nie je v notovom zápise výslovne uvedená, pretože je súčasťou vecí, ako je taktové označenie, ktoré sa objavuje iba raz na začiatku}} (existuje príliš veľa takýchto „vecí`` na to, aby sme ich tu vymenovali, ako napríklad rozdiel medzi valčíkom a mazurkou. Ďalší príklad: bez toho, aby sa pozreli na hudbu, by si niekto myslel, že rytmus v piesni „happy birthday`` je na „happy``, ale v skutočnosti je na „birth``; táto pieseň je valčík). V mnohých prípadoch je hudba vytvorená hlavne manipuláciou s týmito rytmickými variáciami, takže rytmus je jedným z najdôležitejších prvkov hudby. Stručne povedané, väčšina rytmických ťažkostí vzniká z nesprávneho čítania hudby. Toto sa často stáva, keď sa snažíte čítať hudbu HT; je tam jednoducho príliš veľa informácií na spracovanie a mozog sa nemôže starať o rytmus, najmä ak hudba zahŕňa nové technické zručnosti. Táto počiatočná chyba pri čítaní sa potom opakovaným cvičením začlení do výslednej hudby.

\textbf{Definícia rytmu:} Rytmus sa skladá z dvoch častí: načasovania a akcentov a vyskytuje sa v dvoch formách, formálnej a logickej. Záhady okolo rytmu a ťažkosti, s ktorými sa stretávame pri jeho definovaní, vyplývajú z „logickej`` časti, ktorá je zároveň kľúčovým a najťažším prvkom. Najprv sa teda pozrime na jednoduchšie formálne rytmy. Len preto, že sú jednoduchšie, neznamená to, že nie sú dôležité; príliš veľa študentov robí s týmito prvkami chyby, ktoré môžu spôsobiť, že hudba bude nerozpoznateľná.

\textbf{Formálne časovanie: \emph{Formálne časovanie rytmu je dané taktovým označením}} a je uvedené na samom začiatku notového zápisu. Hlavné taktové označenia sú valčík (3/4), bežný takt (4/4), „krátky takt`` (2/2, tiež alla breve) a 2/4. Valčík má 3 doby na takt atď.; počet dôb na takt je označený čitateľom. 4/4 je najbežnejšie a často sa ani neuvádza, hoci by malo byť označené „C`` na začiatku (pamätáte si ho ako „C`` znamená bežný takt). Krátky takt je označený rovnakým „C``, ale so zvislou čiarou v strede (skračuje „C`` na polovicu). Referenčná nota je označená menovateľom, takže 3/4 valčík má 3 štvrťnoty na takt a 2/4 je v zásade dvakrát rýchlejšia ako 2/2 krátky takt. Metrum je počet dôb v takte a takmer každé metrum je zostavené z dvojdôb alebo trojdôb, hoci sa pre špeciálne efekty používajú výnimky (5, 7 alebo 9 dôb).

\textbf{Formálne prízvuky:} Každý taktové označenie má svoj vlastný formálny prízvuk (hlasnejšie alebo tichšie doby). Ak použijeme konvenciu, že 1 je najhlasnejšia, 2 je tichšia atď., potom má (viedenský) valčík formálny prízvuk 133 (známy \textbf{um}-pa-pa); prízvuk dostáva prvá doba; mazurka môže byť 313 alebo 331. Bežný taktový takt má formálny prízvuk 1323 a krátke taktové takty a 2/4 majú prízvuk 1212. Synkopa je rytmus, v ktorom je prízvuk umiestnený na určitom mieste odlišný od formálneho prízvuku; napríklad synkopovaná 4/4 môže byť 2313 alebo 2331. Všimnite si, že rytmus 2331 je v celej skladbe nemenný, ale 1 je na netradičnom mieste.

\textbf{Logické načasovanie a akcenty:} Tu skladateľ vnáša svoju hudbu. Ide o zmenu načasovania a hlasitosti oproti formálnemu rytmu. Hoci rytmická logika nie je nevyhnutná, je takmer vždy prítomná. Bežnými príkladmi rytmickej logiky načasovania sú zrýchlenie (aby boli veci vzrušujúcejšie), spomalenie (možno na označenie konca) alebo rubato. Príkladmi dynamickej rytmickej logiky sú zvyšovanie alebo znižovanie hlasitosti, forte, PP atď.

\textbf{Beethovenova Búrlivá sonáta} (Op. 31, č. 2) ilustruje formálne a logické rytmy. Napríklad v 3. vete sú prvé 3 takty 3 opakovaniami rovnakej štruktúry a jednoducho sledujú formálny rytmus. Avšak v taktoch 43 -- 46 je v pravej note 6 opakovaní rovnakej štruktúry, ale musia byť vtesnané do 4 formálnych rytmických taktov! Ak v pravej note urobíte 6 rovnakých opakovaní, mýlite sa! Okrem toho sa v takte 47 nachádza neočakávané „sf``, ktoré nemá nič spoločné s formálnym rytmom, ale je absolútne nevyhnutným logickým rytmom.

Ak je rytmus taký dôležitý, aké rady možno použiť na jeho pestovanie? Je zrejmé, že \emph{\textbf{rytmus musíte brať ako samostatný predmet cvičenia, pre ktorý potrebujete špecifický program ako na to}}. Preto si počas úvodného učenia sa novej skladby vyhraďte nejaký čas na prácu s rytmom. Metronóm, najmä taký s pokročilými funkciami, môže byť v tomto prípade užitočný. Najprv musíte dvakrát skontrolovať, či je váš rytmus v súlade s taktom. Toto sa nedá urobiť v mysli ani po tom, čo skladbu zahráte -- musíte si znova pozrieť notový zápis a skontrolovať každú notu. Príliš veľa študentov hrá skladbu len určitým spôsobom, „pretože znie správne``; toto nemôžete robiť. Musíte si overiť v partitúre, či správne noty nesú správny prízvuk presne podľa taktového označenia. Až potom sa môžete rozhodnúť, ktorá rytmická interpretácia je najlepší spôsob hrania a kde skladateľ vložil porušenia základných pravidiel (veľmi zriedkavé); častejšie je rytmus naznačený taktom presne správny, ale znie protiintuitívne. Príkladom toho je tajomné „arpeggio`` na začiatku \textbf{Beethovenovej Appassionaty} (Op. 57). Normálne arpeggio (ako napríklad CEG) by začínalo prvou notou (C), ktorá by mala niesť prízvuk (slabšia doba). Beethoven však začína každý takt na tretej note arpeggia (prvý takt je neúplný a nesie prvé dve noty „arpeggia``); to vás núti zdôrazniť tretiu notu (G), nie prvú notu, ak správne dodržiavate taktové označenie. Dôvod tohto zvláštneho „arpeggia`` zistíme, keď je hlavná téma uvedená v takte 35. Všimnite si, že toto „arpeggio`` je len obrátená, schematická (zjednodušená) forma hlavnej témy. Beethoven nás psychologicky pripravil na hlavnú tému tým, že nám dal iba rytmus! Preto ju opakuje po tom, čo ju zvýšil o zvláštny interval -- chcel sa len uistiť, že rozoznáme nezvyčajný rytmus (rovnaký prostriedok použil na začiatku svojej 5. symfónie, kde zopakoval 4-tónový motív v nižšej výške). Ďalším príkladom je Chopinova Fantaisie-Impromptu. Prvá nota RH (takt 5) musí byť tichšia ako druhá. Viete nájsť aspoň jeden dôvod, prečo? Hoci je táto skladba v dvojnásobnom takte, môže byť poučné precvičiť si pravú notu v 4/4, aby ste sa uistili, že nebudú zdôraznené nesprávne tóny.

\emph{\textbf{Keď začínate s hraním rukami zvlášť (HS), pozorne si skontrolujte rytmus. Potom ho skontrolujte znova, keď začínate s hraním rukami spolu (HT). Keď je rytmus nesprávny, hudbu zvyčajne nemožno hrať v rýchlosti. Preto, ak máte nezvyčajné ťažkosti s dostávaním sa do tempa, je dobré rytmus skontrolovať. V skutočnosti je nesprávna rytmická interpretácia jednou z najčastejších príčin rýchlostných bariér a dôvodom, prečo máte problémy s hraním rukami spolu. Keď urobíte rytmickú chybu, žiadne množstvo cvičenia vám nepomôže dosiahnuť tempo!}} Toto je jeden z dôvodov, prečo osnovovanie funguje: môže zjednodušiť úlohu správneho čítania rytmu. Preto pri osnovovaní sa sústreďte na rytmus. Taktiež, keď s HT začínate, môžete mať väčší úspech, ak rytmus preženiete. Rytmus je ďalším dôvodom, prečo by ste sa nemali pokúšať o skladby, ktoré sú pre vás príliš ťažké. Ak nemáte dostatočnú techniku, nebudete schopní ovládať rytmus. Môže sa stať, že nedostatok techniky vnúti vašej hre nesprávny rytmus, čím sa vytvorí rýchlostná stena.

Ďalej hľadajte špeciálne rytmické označenia, ako napríklad „s\textbf{f}” alebo prízvuky. Nakoniec existujú situácie, v ktorých na hudbe nie sú žiadne označenia a vy jednoducho musíte vedieť, čo robiť, alebo si vypočuť nahrávku, aby ste zachytili špeciálne rytmické variácie. Preto by ste ako súčasť cvičenia mali experimentovať s rytmom, zdôrazňovať neočakávané tóny atď., aby ste zistili, čo sa môže stať.

Rytmus je tiež úzko spojený s rýchlosťou. Preto je potrebné hrať väčšinu Beethovenových skladieb nad určitými rýchlosťami; inak sa môžu stratiť emócie spojené s rytmom a dokonca aj melodické linky. Beethoven bol majstrom rytmu; preto Beethovena nemôžete úspešne hrať bez toho, aby ste venovali osobitnú pozornosť rytmu. Zvyčajne vám dáva aspoň dve veci súčasne: (i) ľahko sledovateľnú melódiu, ktorú publikum počuje, a (ii) rytmický/harmonický prostriedok, ktorý ovláda to, čo publikum \emph{cíti}. Tak v prvej vete jeho Patetickej (Op. 13) rozrušený tremolo LH ovláda emócie, zatiaľ čo publikum je zaneprázdnené počúvaním zvedavého pravého rytmu. Preto len technická schopnosť zvládnuť rýchle tremolo LH nestačí -- musíte byť schopní ovládať emocionálny obsah pomocou tohto tremola. Keď pochopíte a dokážete vykonať rytmický koncept, bude oveľa jednoduchšie vyzdvihnúť hudobný obsah celej vety a ostrý kontrast s časťou Grave sa stane zrejmým.

Existuje jedna trieda rytmických ťažkostí, ktoré sa dajú vyriešiť jednoduchým trikom. Je to trieda zložitých rytmov s chýbajúcimi notami. Dobrým príkladom je 2. veta Beethovenovej Patetickej. Rytmické označenie 2/4 sa v taktoch 17 až 21 ľahko hrá kvôli opakujúcim sa akordom ľavej ruky (LH), ktoré udržiavajú rytmus. V takte 22 však chýbajú najdôležitejšie akcentované noty ľavej ruky, čo sťažuje zachytenie trochu zložitej hry v pravej ruke. Riešením tohto problému je jednoducho doplniť chýbajúce noty ľavej ruky! Týmto spôsobom si môžete ľahko precvičiť správny rytmus v pravej ruke.

Stručne povedané, „tajomstvo`` skvelého rytmu nie je žiadnym tajomstvom -- musí začať správnym počítaním (čo, ako musím znova zdôrazniť, nie je jednoduché). Pre pokročilých klaviristov je to samozrejme oveľa viac; je to mágia. To je to, čo odlišuje skvelého od obyčajného. Nejde len o počítanie prízvukov v každom takte, ale aj o to, ako sa takty spájajú a vytvárajú rozvíjajúcu sa hudobnú myšlienku -- logickú zložku rytmu. Napríklad v Beethovenovom Mesačnom svite (Op. 27) je začiatok 3. vety v podstate 1. veta hraná šialenou rýchlosťou. Táto znalosť nám hovorí, ako hrať 1. vetu, pretože to znamená, že séria triol v 1. vete musí byť prepojená takým spôsobom, aby viedla k vyvrcholeniu tromi opakovanými notami. Ak by ste opakované noty hrali jednoducho nezávisle od predchádzajúcich triol, všetky tieto noty stratia svoj význam/dopad. Rytmus je tiež zvláštny alebo neočakávaný prízvuk, ktorý náš mozog nejako rozpoznáva ako zvláštny. Rytmus je jednoznačne kritickým prvkom hudby, ktorému musíme venovať osobitnú pozornosť.

\begin{enumerate}
\def\labelenumi{\alph{enumi}.}
\setcounter{enumi}{2}
\item
  \textbf{Legato, Staccato}
\end{enumerate}

Legato je plynulá hra. To sa dosahuje prepojením po sebe nasledujúcich nôt -- prvú notu nezdvíhajte, kým nezahráte druhú. Fraser odporúča značné prekrývanie dvoch nôt. Prvé momenty noty obsahujú veľa „šumu``, takže prekrývajúce sa noty nie sú až také viditeľné. Keďže legato je zvyk, ktorý si musíte do hry zabudovať, experimentujte s rôznymi úrovňami prekrývania, aby ste zistili, aké prekrývanie \emph{vám} dáva najlepšie legato. Potom to precvičujte, kým sa z toho nestane zvyk, aby ste vždy dokázali reprodukovať ten istý efekt. Chopin považoval legato za najdôležitejšiu zručnosť, ktorú si má začiatočník osvojiť.

Chopinova hudba vyžaduje špeciálne druhy legata a staccata (Balada op. 23), preto je dôležité venovať týmto prvkom pozornosť pri hraní jeho hudby. \emph{\textbf{Ak si chcete precvičovať legato, zahrajte si Chopina.}} Základný klávesový úder je pre legato absolútne nevyhnutný.

\emph{\textbf{Pri staccate sa prst odrazí od klávesu, čím sa vytvorí krátky zvuk bez doznievania.}} Je trochu prekvapujúce, že väčšina kníh o učení sa hry na klavíri sa zaoberá staccatom, ale nikdy nedefinuje, čo to je! Pri staccate nie je aktivovaný spätný ráz a tlmič preruší zvuk ihneď po zahraní noty. Preto chýba „držiaca`` zložka základného úderu. Existujú dva zápisy pre staccato, normálny (bodka) a tvrdý (plný trojuholník). V oboch prípadoch sa zdvihák neuvoľní; pri tvrdom staccate sa prst pohybuje nadol a nahor oveľa rýchlejšie. Pri normálnom staccate teda môže byť stlačenie klávesu približne v polovici, ale pri tvrdom staccate to môže byť menej ako v polovici. Týmto spôsobom sa tlmič vracia na struny rýchlejšie, čo má za následok kratšie trvanie noty. Pretože spätný ráz nie je aktivovaný, kladivko sa môže „odrážať``, čo sťažuje opakovanie pri určitých rýchlostiach. Ak teda máte problémy s rýchlo sa opakujúcimi staccatmi, neobviňujte sa hneď -- môže to byť nesprávna frekvencia, pri ktorej kladivko odráža nesprávnym smerom. Zmenou rýchlosti, miery poklesu kláves atď. môžete tento problém odstrániť. Pri normálnom staccate gravitácia rýchlo vracia tlmič na struny. Pri tvrdom staccate sa tlmič v skutočnosti odráža od hornej lišty tlmiča, takže sa vracia ešte rýchlejšie. Ohyb drieku kladivka môže byť pri staccate negatívny, čo znižuje efektívnu hmotnosť kladivka; preto existuje značná škála tónov, ktoré môžete so staccatom vytvoriť. Preto sa pri staccate menia pohyby kladivka, spätný tlmič, zdviháka a tlmiča. Je \emph{\textbf{zrejmé, že na dobrú hru staccato je užitočné pochopiť, ako klavír funguje.}}

Staccato sa vo všeobecnosti delí na tri typy v závislosti od spôsobu hry: (i) staccato prstami, (ii) staccato zápästím a (iii) staccato rukou, ktoré zahŕňa pohyb hore a dole aj rotáciu ruky. (i) hrá sa prevažne prstami, pričom ruka a rameno zostávajú nehybné, (ii) hrá sa prevažne zápästím a (iii) zvyčajne sa hrá ako ťah (III.4a), pričom hra začína v hornej časti paže. Ako postupujete od (i) k (iii), pridávate viac hmoty za prstami; preto (i) dáva najľahšie a najrýchlejšie staccato a je užitočné pre jednotlivé, tiché tóny a (iii) dáva najťažší pocit a je užitočné pre hlasné pasáže a akordy s mnohými tónmi, ale je tiež najpomalšie. (ii) je medzi tým. V praxi väčšina z nás pravdepodobne kombinuje všetky tri; keďže zápästie a ruka sú pomalšie, ich amplitúdy sa musia zodpovedajúcim spôsobom znížiť, aby sa dalo hrať rýchle staccato. Niektorí učitelia sa mračia na používanie staccato zápästím a uprednostňujú prevažne staccato rukou; Pravdepodobne je však lepšie mať na výber (alebo kombináciu) všetkých troch. Napríklad, únavu by ste mohli znížiť zmenou jedného za druhý, hoci štandardnou metódou znižovania únavy je striedanie prstov. Pri cvičení staccata si najprv precvičte tri staccata (prsty, ruka, paža) a až potom sa rozhodnete, ktoré použijete alebo ako ich skombinujete.

Keďže pri staccato nemôžete použiť váhu paží, najlepšou referenciou je vaše stabilné telo. Telo teda zohráva pri staccato hre dôležitú úlohu. Rýchlosť opakovania staccata je riadená množstvom pohybu hore a dole: čím menší pohyb, tým rýchlejšia je frekvencia opakovania, presne rovnako ako pri driblingu basketbalovej lopty.

\begin{enumerate}
\def\labelenumi{\arabic{enumi}.}
\setcounter{enumi}{1}
\item
  \textbf{Cyklovanie} (Chopinova Fantaisie Impromptu)
\end{enumerate}

\emph{\textbf{Cyklovanie je najlepší postup na budovanie techniky pre veci, ako sú nové alebo rýchle pasáže, ktoré nezvládate. Cyklovanie (tiež nazývané „opakovanie dookola``) je vziať segment a hrať ho opakovane; zvyčajne nepretržite, bez prestávok.}} Ak je spojka potrebná na nepretržité cyklovanie rovnaká ako prvá nota segmentu, tento segment cykluje prirodzene; nazýva sa to samocyklujúci segment. Príkladom je CGEG štvorica. Ak je spojka iná, musíte si vymyslieť takú, ktorá vedie k prvej note, aby ste mohli cyklovať bez prestávok.

\emph{\textbf{Cyklovanie je v podstate čisté opakovanie, ale je dôležité používať ho takmer ako postup proti opakovaniu, spôsob, ako sa vyhnúť bezmyšlienkovému opakovaniu. Myšlienkou cyklovania je, že si techniku osvojíte tak rýchlo, že eliminujete zbytočné, bezmyšlienkovité opakovanie.}} Aby ste sa vyhli osvojovaniu si zlých návykov, zmeňte rýchlosť a experimentujte s rôznymi polohami rúk/paží/prstov pre optimálnu hru a vždy sa snažte relaxovať; snažte sa necyklovať presne to isté príliš veľakrát. Hrajte potichu (aj hlasné úseky), kým nezvládnete techniku, dosahujte rýchlosť aspoň o 20 \% vyššiu ako konečná rýchlosť a ak je to možné, až dvojnásobok konečnej rýchlosti. Viac ako 90 \% času cyklovania by malo byť pri rýchlostiach, ktoré zvládnete pohodlne a presne. Potom postupne znižujte rýchlosť na veľmi pomalú. Hotovo je, keď dokážete hrať akoukoľvek rýchlosťou po akúkoľvek dobu, bez toho, aby ste sa pozerali na ruku, úplne uvoľnene a s plnou kontrolou. Možno zistíte, že určité stredné rýchlosti spôsobujú problémy. Precvičujte si tieto rýchlosti, pretože môžu byť potrebné na začiatku HT. Cvičte bez pedálu (čiastočne preto, aby ste sa vyhli zlozvyku nestláčať úplne cez klávesový úder), kým si techniku nezvládnete. \emph{\textbf{Často striedajte ruky, aby ste predišli zraneniu.}}

Ak technika vyžaduje 10 000 opakovaní (typická požiadavka pre naozaj náročný materiál), cyklovanie vám umožní ich zvládnuť v čo najkratšom čase. Reprezentatívne časy cyklov sú približne 1 sekunda, takže 10 000 cyklov je menej ako 4 hodiny. Ak budete tento segment cyklovať 10 minút denne, 5 dní v týždni, 10 000 cyklov vám zaberie takmer mesiac. Je zrejmé, že naučenie sa veľmi náročného materiálu bude trvať mesiace s použitím najlepších metód a \emph{oveľa} dlhšie, ak použijete menej efektívne metódy.

\emph{\textbf{Cyklovanie je potenciálne najškodlivejším zo všetkých cvičení na klavíri}}, preto buďte opatrní. Nepreháňajte to prvý deň a uvidíte, čo sa stane na druhý deň. Ak vás na druhý deň nič nebolí, môžete v tréningu cyklovania pokračovať alebo ho predĺžiť. Predovšetkým, vždy pri cyklovaní cvičte na dvoch cykloch naraz, jednom pre pravú a druhom pre ľavú ruku, aby ste si mohli často striedať ruky. U mladých ľudí môže nadmerné cyklovanie spôsobiť bolesť; v takom prípade s cyklovaním prestaňte a ruka by sa mala zotaviť o niekoľko dní. U starších ľudí môže nadmerné cyklovanie spôsobiť osteoartritické vzplanutia, ktoré môžu ustúpiť až po niekoľkých mesiacoch.

Aplikujme cyklovanie na Chopinovu FI: ľavoručné arpeggio, takt 5. Prvých šesť nôt cykluje samostatne, takže to môžete skúsiť. Keď som to skúsil prvýkrát, natiahnutie bolo pre moje malé ruky príliš veľké, takže som sa príliš rýchlo unavil. Urobil som to tak, že som cykloval prvých 12 nôt. Druhých, ľahších šesť nôt mi umožnilo trochu si oddýchnuť v rukách, a preto som mohol cyklovať 12-notový segment dlhšie a rýchlejšie. Samozrejme, ak naozaj chcete zvýšiť rýchlosť (nie je to potrebné pre ľavú ruku, ale v tejto skladbe by to mohlo byť užitočné pre pravú), cyklujte iba prvú paralelnú sadu (prvé tri alebo štyri noty pre ľavú ruku).

Len preto, že viete zahrať prvý segment, neznamená to, že teraz viete zahrať všetky ostatné arpeggia. Budete musieť začať prakticky od nuly, dokonca aj pre rovnaké noty o oktávu nižšie. Samozrejme, druhé arpeggio bude po zvládnutí prvého jednoduchšie, ale možno budete prekvapení, koľko práce musíte zopakovať, keď sa v segmente vykoná veľmi malá zmena. Deje sa to preto, lebo v tele je toľko svalov, že váš mozog si môže vybrať inú sadu svalov na vytvorenie pohybov, ktoré sa len mierne líšia (a zvyčajne to tak aj je). Na rozdiel od robota nemáte veľa možností, ktoré svaly si mozog vyberie. Až keď ste zahrali veľmi veľký počet takýchto arpeggií, ďalšie príde ľahko. Preto by ste mali počítať s tým, že budete musieť cyklicky hrať pomerne málo arpeggií.

Aby sme pochopili, ako hrať túto Chopinovu skladbu, je užitočné analyzovať matematický základ časovania 3 verzus 4 v tejto skladbe. Pravá ruka hrá veľmi rýchlo, povedzme 4 noty za pol sekundy (približne). Zároveň ľavá ruka hrá pomalšie, 3 noty za pol sekundy. Ak sú všetky noty zahrané presne, publikum počuje frekvenciu nôt ekvivalentnú 12 notám za pol sekundy, pretože táto frekvencia zodpovedá najmenšiemu časovému intervalu medzi notami. \emph{\textbf{To znamená, že ak vaša pravá ruka hrá tak rýchlo, ako len dokáže, potom pridaním POMALŠEJ hry s ľavou rukou sa Chopinovi podarilo zrýchliť túto skladbu na 3-násobok vašej maximálnej rýchlosti}} !

Ale počkajte, nie je tam všetkých 12 nôt; v skutočnosti ich je len 7, takže chýba 5 nôt.

Tieto chýbajúce tóny vytvárajú to, čo sa nazýva Moiré vzor, čo je tretí vzor, ktorý vzniká superpozíciou dvoch nesúmerných vzorov. Tento vzor vytvára v každom takte vlnový efekt a Chopin tento efekt zosilnil použitím arpeggia LH, ktoré stúpa a klesá ako vlna synchrónne s Moiré vzorom. Zrýchlenie faktorom 3 a Moiré vzor sú záhadné efekty, ktoré sa publiku nenápadne vynárajú, pretože netušia, čo ich vytvorilo, alebo že vôbec existujú. Mechanizmy, ktoré ovplyvňujú publikum bez jeho vedomia, často vytvárajú dramatickejšie efekty ako tie, ktoré sú zrejmé (ako napríklad hlasné hranie, legato alebo rubato). Veľkí skladatelia vynašli neuveriteľné množstvo týchto skrytých mechanizmov a matematická analýza je často najjednoduchší spôsob, ako ich odhaliť. Chopin pravdepodobne nikdy neuvažoval v pojmoch nesúmerných množín a Moiré vzorov; tieto koncepty intuitívne chápal vďaka svojmu géniu.

Je poučné špekulovať o dôvode chýbajúcej prvej noty taktu (5. takt) v pravej ruke, pretože ak dokážeme dôvod rozlúštiť, budeme presne vedieť, ako to zahrať. Všimnite si, že sa to deje na samom začiatku melódie pravej ruky. Na začiatku melódie alebo hudobnej frázy sa skladatelia vždy stretávajú s dvoma protichodnými požiadavkami: jednou je, že každá fráza by mala (vo všeobecnosti) začínať jemne, a druhou je, že prvá nota taktu je pomalá a mala by byť prízvukovaná. Skladateľ môže úhľadne splniť obe požiadavky odstránením prvej noty, čím zachová rytmus a napriek tomu začne jemne (v tomto prípade žiadny zvuk)! Nebudete mať problém nájsť množstvo príkladov tohto prostriedku -- pozri Bachove invencie. Ďalším prostriedkom je začať frázu na konci čiastočného taktu tak, aby prvá pomalá doba prvého celého taktu prišla po odohraní niekoľkých nôt (klasickým príkladom je začiatok prvej vety Beethovenovej Appassionaty). To znamená, že prvá nota pravej ruky v tomto takte Chopinovho FI musí byť tichá a druhá nota hlasnejšia ako prvá, aby sa presne zachoval rytmus (ďalší príklad dôležitosti rytmu!). Nie sme zvyknutí hrať týmto spôsobom; bežná hra je začať prvú notu ako ťažkú dobu. V tomto prípade je to obzvlášť ťažké kvôli rýchlosti; preto si tento začiatok môže vyžadovať dodatočný začiatok.

Táto skladba začína postupným vťahovaním publika do svojho rytmu ako neodolateľné pozvanie, po tom, čo na seba upozorní hlasnou oktávou prvého taktu, po ktorej nasleduje rytmické arpeggio v spodnej osnove. Chýbajúca nota v piatom takte je po niekoľkých opakovaniach obnovená, čím sa zdvojnásobí frekvencia opakovania Moiré a efektívny rytmus. V druhej téme (13. takt) je plynulá melódia pravej noty nahradená dvoma prerušenými akordmi, čím sa vytvára dojem štvornásobného zosilnenia rytmu. Toto „rytmické zrýchlenie`` vrcholí vo vrcholnom forte taktov 19 -- 20. Publikum si potom oddýchne „zjemnením`` rytmu vytvoreným oneskoreným melodickým (malíčkovým) tónom pravej ruky a potom jeho postupným doznievaním, ktoré sa dosiahne diminuendom až po PP. Celý cyklus sa potom opakuje, tentoraz s pridanými prvkami, ktoré zvyšujú vyvrcholenie, až kým nekončí v zostupných, dunivých rozložených akordoch. Na precvičovanie tejto časti je možné každý rozložený akord cyklovať samostatne. Tieto akordy neobsahujú konštrukt 3,4 a vynášajú vás späť z tajomného podsvetia 3,4, čím vás pripravujú na pomalú časť.

Rovnako ako u väčšiny Chopinových skladieb, ani pre túto skladbu neexistuje „správne`` tempo. Ak však hráte rýchlejšie ako približne 2 sekundy/takt, efekt násobenia 3x4 má tendenciu miznúť a zvyčajne vám zostane prevažne efekt Moiré a ďalšie efekty. Je to čiastočne kvôli klesajúcej presnosti s rýchlosťou, ale čo je dôležitejšie, pretože rýchlosť 12x sa stáva príliš rýchlou na to, aby ju ucho dokázalo sledovať. Nad približne 20 Hz začínajú opakovania nadobúdať pre ľudské ucho vlastnosti zvuku. Preto multiplikačné zariadenie funguje iba do približne 20 Hz; nad touto frekvenciou získate nový efekt, ktorý je ešte výnimočnejší ako neuveriteľná rýchlosť -- „rýchle tóny`` sa menia na „nízkofrekvenčný zvuk``. 20 Hz je teda akýsi zvukový prah. Preto je najnižšia nota klavíra A s frekvenciou približne 27 Hz. Tu je veľké prekvapenie: existujú dôkazy, že Chopin tento efekt počul! Všimnite si, že rýchla časť je spočiatku označená ako „Allegro agitato``, čo znamená, že každá nota musí byť jasne počuteľná. Na metronóme zodpovedá Allegro 12-násobnej rýchlosti od 10 do 20 Hz, čo je správna frekvencia na počutie... násobenie, tesne pod „zvukovým prahom``. „Agitato`` zabezpečuje, že táto frekvencia je počuteľná. Keď sa táto rýchla časť vráti po časti Moderato, je označená ako Presto, čo zodpovedá 20 až 40 Hz -- chcel, aby sme ju hrali pod a nad zvukovým prahom! Preto existujú matematické dôkazy naznačujúce, že Chopin o tomto zvukovom prahu vedel.

Pomalá stredná časť bola stručne opísaná v časti II.25. Najrýchlejší spôsob, ako sa ju naučiť, podobne ako mnoho Chopinových skladieb, je začať memorovaním ľavej ruky (LH). Je to preto, že postup akordov často zostáva rovnaký, aj keď Chopin nahradí pravý akord novou melódiou, pretože ľavá ruka poskytuje hlavne sprievodné akordy. Všimnite si, že časovanie 4,3 je teraz nahradené časovaním 2,3 hraným oveľa pomalšie. Používa sa na iný efekt, na zjemnenie hudby a umožnenie voľnejšieho, tempového rubata.

Tretia časť je podobná prvej, až na to, že sa hrá rýchlejšie, čo má za následok úplne iný efekt, a záver je odlišný. Tento záver je náročný pre malé ruky a môže vyžadovať extra cyklovanie pravej ruky. V tejto časti nesie melódiu pravý malíček, ale zodpovedajúca oktávová nota hraná palcom obohacuje melodickú linku. Skladba končí nostalgickým preformulovaním témy pomalého pohybu v ľavej ruke. Vrchnú notu tejto melódie v ľavej ruke (G\# - takt 7 od konca) jasne odlíšte od tej istej noty, ktorú hrá pravá ruka, tak, že ju podržíte o niečo dlhšie a potom ju podržíte pedálom.

G\# je najdôležitejšia nota v tejto skladbe. Začiatočná oktáva s\emph{\textbf{f}} G\# teda nie sú len fanfáry uvádzajúce skladbu, ale aj šikovný spôsob, akým Chopin vštepil G\# do myslí poslucháčov. Preto sa s touto notou neponáhľajte; dajte si načas a nechajte ju vstrebať. Ak si prečítate celý tento text, uvidíte, že G\# zaberá všetky dôležité pozície. V pomalej časti je G\# A♭, čo je tá istá nota. Tento G\# je ďalší z tých nástrojov, v ktorých skvelý skladateľ opakovane „udieranie publika po hlave pálkou dva krát štyri`` (G\#), ale publikum o tom nemá ani tušenie čo ich zasiahlo. Pre klaviristu znalosť G\# pomáha interpretovať a zapamätať si skladbu. Teda koncepčné vyvrcholenie tejto skladby prichádza na konci (ako by aj malo byť), keď musia obe ruky hrať rovnaké G\# (takty 8 a 7 od konca); preto sa toto ľavo-pravé G\# musí vykonať s maximálnou starostlivosťou, pričom sa zachováva plynule slabnúca oktáva RH G\#.

Naša analýza jasne zdôrazňuje otázku, ako rýchlo hrať túto skladbu. Na dosiahnutie efektu dvanásťnôtového rytmu a neľudsky presného hrania nad zvukovým prahom je potrebná vysoká presnosť. Ak sa túto skladbu učíte prvýkrát, frekvencia dvanástich nôt nemusí byť spočiatku počuteľná kvôli nedostatočnej presnosti. Keď to konečne „pochopíte``, hudba bude zrazu znieť veľmi „rušne``. Ak hráte príliš rýchlo a stratíte presnosť, môžete stratiť faktor tri -- jednoducho sa vymyje a publikum počuje iba štyri noty. Pre začiatočníkov sa dá skladba zrýchliť spomalením a zvýšením presnosti. Hoci pravá ruka nesie melódiu, ľavá ruka musí byť jasne počuteľná; inak zmizne efekt dvanástich nôt aj moiré. Keďže ide o Chopinovu skladbu, nie je potrebné, aby bol počuť efekt dvanástich nôt; táto skladba je prístupná nekonečnému množstvu interpretácií a niektorí možno budú chcieť potlačiť ľavú ruku a sústrediť sa na pravú ruku a stále vytvoriť niečo magické.

Výhodou cyklovania je, že ruka hrá nepretržite, čo simuluje nepretržité hranie lepšie, ako keby ste cvičili izolované segmenty. Umožňuje vám to tiež experimentovať s malými zmenami v polohe prstov atď., aby ste našli optimálne podmienky pre hranie. Nevýhodou je, že pohyby rúk pri cyklovaní sa môžu líšiť od pohybov potrebných na hranie skladby. Ramená majú pri cyklovaní tendenciu byť nehybné, zatiaľ čo v skutočnej skladbe sa ruky zvyčajne musia pohybovať. Preto v prípadoch, keď segment prirodzene necykluje, možno budete musieť použiť segmentálny tréning bez cyklovania. Jednou z výhod necyklovania je, že teraz môžete zahrnúť spojku.

\textbf{3. Trilky a tremolá}

\begin{enumerate}
\def\labelenumi{\alph{enumi}.}
\item
  \textbf{Trilky}
\end{enumerate}

\emph{\textbf{Neexistuje lepšia demonštrácia účinnosti cvičení s paralelnými sadami (PS) (pozri III.7) ako ich použitie na učenie sa trilku. Pri trilkovaní je potrebné vyriešiť dva hlavné problémy: (1) rýchlosť (s kontrolou) a (2) pokračovať v hraní tak dlho, ako je potrebné.}} Cvičenia s PS boli navrhnuté tak, aby riešili presne tieto typy problémov, a preto fungujú veľmi dobre na nácvik trilkov. Whiteside opisuje metódu nácviku trilku, ktorá je typom akordického útoku. Použitie akordického útoku na nácvik trilku teda nie je ničím novým. Keďže však teraz mechanizmus učenia chápeme podrobnejšie, môžeme navrhnúť najpriamejší a najefektívnejší prístup pomocou PS.

Prvým problémom, ktorý treba vyriešiť, sú prvé dve noty. \emph{\textbf{Ak sa prvé dve noty nezačnú správne hrať, naučiť sa trilkovať sa stáva náročnou úlohou. Dôležitosť prvých dvoch nôt platí aj pre behy, arpeggia atď.}} Riešenie je však takmer triviálne -- použite cvičenie PS s dvoma notami. Preto pre trilok 2323.... použite prvú trojku ako spojku a precvičte si 23. Potom si precvičte 32, potom 232 atď. Je to také jednoduché! Vyskúšajte to! Funguje to ako kúzlo! Predtým, ako aplikujete PS na trilok, si možno budete chcieť prečítať časť III.7 o nich.

\emph{\textbf{Trilok pozostáva z 2 pohybov: pohybu prstov a rotácie predlaktia. Preto si tieto 2 zručnosti precvičujte samostatne.}} Najprv používajte na trilkovanie iba prsty, pričom ruka a rameno sú úplne nehybné. Potom držte prsty fixované a trilkujte iba s rotáciou ramena. Takto zistíte, či vás spomaľujú prsty alebo rotácia ramena. Mnoho študentov nikdy necvičilo rýchlu rotáciu ramena (kývanie ramena) a často to bude pomalší pohyb. Pri rýchlych trilkoch je táto rotácia tam a späť neviditeľne malá, ale nevyhnutná. Aplikujte cvičenia PS na rotačné pohyby prstov aj ramena. Pri pomalých trilkoch pohyby preháňajte a rýchlosť zvyšujte znížením veľkosti pohybov. Konečná veľkosť oboch pohybov nemusí byť rovnaká, pretože pri pomalšom pohybe (rotácia ramena) použijete menší pohyb, aby ste kompenzovali jeho pomalosť. Pri precvičovaní týchto pohybov experimentujte s rôznymi polohami prstov. Pozrite si časť Tremolo, kde platia podobné metódy -- trilok je len zmenšené tremolo.

\emph{\textbf{Relaxácia je pre trilok ešte dôležitejšia ako takmer akákoľvek iná technika}} kvôli potrebe rýchleho vyváženia hybnosti; to znamená, že PS sú len dva tóny, existuje príliš veľa spojok na to, aby sme sa spoliehali výlučne na paralelizmus na dosiahnutie rýchlosti. Preto musíme byť schopní rýchlo meniť hybnosť prstov. Pri trilkoch musí byť hybnosť pohybu prstov vyvážená rotáciou paží. Napätie uzamkne prsty k väčším častiam, ako sú dlane a ruky, čím sa zvýši efektívna hmotnosť prstov. Väčšia hmotnosť znamená pomalší pohyb: pozrite sa na fakt, že kolibrík dokáže mávať krídlami rýchlejšie ako kondor a malý hmyz ešte rýchlejšie ako kolibrík. Platí to aj v prípade, že by sa odpor vzduchu ignoroval; v skutočnosti je vzduch pre kolibríka efektívne viskóznejší ako pre kondora a pre malý hmyz je vzduch takmer rovnako viskózny ako voda pre veľkú rybu; napriek tomu hmyz dokáže mávať krídlami rýchlo, pretože hmotnosť krídel je taká malá. Preto je dôležité od samého začiatku začleniť do trilku úplnú relaxáciu, čím sa uvoľnia prsty od ruky. \emph{\textbf{Trilkovanie je jedna zručnosť, ktorá si vyžaduje neustále udržiavanie. Ak sa chcete stať dobrým trilkovačom, budete musieť trilkovať každý deň.}} PS: Cvičenie č. 1. (2-nota) je najlepší postup na udržanie trilku v perfektnom stave, najmä ak ste ho nejaký čas nepoužívali alebo ak ho chcete naďalej vylepšovať.

Nakoniec, trilok nie je séria staccat. Končeky prstov musia byť čo najdlhšie v spodnej časti klávesy, t. j. chytače musia byť zapojené pre každú notu. Dávajte si pozor na minimálny zdvih potrebný na to, aby opakovanie fungovalo. Tí, ktorí zvyčajne cvičia na koncertnom krídle si treba uvedomiť, že táto vzdialenosť zdvihu môže byť pri pianíne takmer dvojnásobná. Rýchlejšie trilky vyžadujú menšie zdvihy, preto na pianíne možno budete musieť trilok spomaliť. Rýchle trilky na elektronických klavíroch sú náročné, pretože ich mechanika je horšia.

\begin{enumerate}
\def\labelenumi{\alph{enumi}.}
\setcounter{enumi}{1}
\item
  \textbf{Tremolá} (Beethovenova patetická, 1. časť)
\end{enumerate}

Tremolá sa cvičia presne rovnakým spôsobom ako trilky. Aplikujme to na niekedy obávané dlhé oktávové tremolá z Beethovenovej Patetickej sonáty (Opus 13). Pre niektorých študentov sa tieto tremolá zdajú nemožné a mnohí si nadmerným cvičením poranili ruky, niektorí natrvalo. Iní s nimi nemajú žiadne ťažkosti. Ak viete, ako ich cvičiť, sú v skutočnosti celkom jednoduché. Posledná vec, ktorú \emph{\textbf{chcete robiť, je cvičiť toto tremolo celé hodiny v nádeji, že si vybudujete vytrvalosť -- to je najistejší spôsob, ako si osvojiť zlé návyky a utrpieť zranenie.}}

Keďže potrebujete oktávové tremolá v oboch rukách, budeme cvičiť ľavú ruku (LH) a striedať sa s cvičením pravej ruky (RH); ak sa pravá ruka uchytí rýchlejšie, môžete ju použiť na výučbu ľavej ruky. Navrhnem postupnosť cvičebných metód; ak máte nejakú predstavivosť, mali by ste byť schopní vytvoriť si vlastnú postupnosť, ktorá by pre vás mohla byť lepšia -- môj návrh je len: návrh na ilustračné účely. Pre úplnosť som to urobil príliš podrobným a príliš dlhým. V závislosti od vašich špecifických potrieb a slabostí by ste mali byť schopní skrátiť postupnosť cvičenia.

Aby ste si precvičili tremolo C2-C3, najprv si precvičte oktávu C2-C3 (ľavá ruka). Pohodlne poskakujte rukou hore a dole a opakujte oktávu s dôrazom na relaxáciu -- dokážete hrať oktávu bez únavy alebo stresu, najmä keď zrýchľujete? Ak sa unavíte, nájdite spôsoby, ako opakovať oktávu bez toho, aby ste sa unavili, zmenou polohy ruky, pohybu atď. Napríklad môžete postupne zdvihnúť zápästie a potom ho znova spustiť -- týmto spôsobom môžete použiť 4 polohy zápästia pre každú štvoricu nôt. Ak ste stále unavení, zastavte sa a vymeňte ruku; precvičte si oktávu RH Ab4-Ab5, ktorú budete neskôr potrebovať. Keď budete môcť hrať opakujúcu sa oktávu 4-krát za dobu (vrátane správneho rytmu) bez únavy, skúste to zrýchliť. Pri maximálnej rýchlosti sa u vás opäť vyvinie únava, preto buď spomaľte, alebo skúste nájsť iné spôsoby hry, ktoré znižujú únavu. Vymeňte ruky hneď, ako sa cítite unavení. \emph{\textbf{Nehrajte nahlas; jedným z trikov na zníženie únavy je cvičiť potichu. Dynamiku môžete pridať neskôr, keď si osvojíte techniku.}} Je mimoriadne dôležité cvičiť potichu. aby ste sa mohli sústrediť na techniku a relaxáciu. Na začiatku, keď sa budete namáhať hrať rýchlejšie, objaví sa únava. Ale keď nájdete správne pohyby, polohy rúk atď., skutočne pocítite, ako únava z ruky odchádza a mali by ste byť schopní oddýchnuť si a dokonca si ruku \emph{pri rýchlom hraní aj osviežiť}. Naučili ste sa relaxovať!

Rovnako ako trilok, aj tremolo pozostáva z pohybu prstov a rotácie rúk. Najprv si precvičte tremolo prstami s použitím prehnaných pohybov prstov, hrajte veľmi pomalé tremolo, zdvíhajte prsty čo najvyššie a silou ich spúšťajte do klávesov. To isté platí pre rotáciu rúk: zafixujte prsty a hrajte tremolo iba s použitím rotácie rúk, prehnane. Všetky pohyby hore a dole musia byť rýchle; ak chcete hrať pomaly, jednoducho medzi pohybmi počkajte a počas tohto čakania precvičujte rýchle a úplné uvoľnenie. Teraz ich postupne zrýchľujte; to sa dosiahne znížením pohybov. Keď je každý z nich uspokojivý, skombinujte ich; pretože oba pohyby prispievajú k tremolu, potrebujete len veľmi málo z každého, a preto budete môcť hrať veľmi rýchlo.

Rýchlosť môžete ešte viac zvýšiť pridaním cvičení PS k cvičeniam na rotáciu prstov aj paží alebo k ich kombinácii. Najprv 5,1 PS. Začnite s opakovanými oktávami a potom postupne každú oktávu nahraďte PS. Napríklad, ak hráte skupiny 4 oktáv (4/4 takt), začnite nahradením 4. oktávy PS, potom 4. a 3. oktávou atď. Čoskoro by ste mali precvičovať všetky PS. Ak sa PS stanú nerovnomernými alebo sa ruka začne unavovať, vráťte sa k oktáve, aby ste si oddýchli. Alebo vymeňte ruky. Pracujte s PS, kým nebudete vedieť zahrať 2 noty v PS takmer „nekonečne rýchlo`` a reprodukovateľne a nakoniec s dobrou kontrolou a úplnou relaxáciou. Pri najrýchlejších rýchlostiach PS by ste mali mať problém rozlišovať medzi PS a oktávami. Potom spomaľte PS tak, aby ste mohli hrať všetkými rýchlosťami s kontrolou. Všimnite si, že v tomto prípade by mala byť nota 5 o niečo hlasnejšia ako 1. Mali by ste to však cvičiť oboma spôsobmi: s dobovým rytmom na 5 aj s ním na 1, aby ste si vytvorili vyváženú a ovládateľnú techniku. Teraz celý postup zopakujte s PS 1,5. Aj keď tento PS nie je potrebný na hranie tohto tremola (potrebný je iba predchádzajúci), je užitočný na rozvoj vyváženej kontroly. Keď sú 5,1 aj 1,5 uspokojivé, prejdite na 5,1,5 alebo 5,1,5,1 (hrané ako krátky oktávový trilok). Ak dokážete 5,1,5,1 zahrať hneď, nie je potrebné robiť 5,1,5. Cieľom je rýchlosť aj vytrvalosť, takže by ste mali cvičiť rýchlosti, ktoré sú \emph{oveľa} rýchlejšie ako konečná rýchlosť tremola, aspoň pri týchto krátkych tremolách. Potom pracujte na 1, 5, 1, 5.

Keď sú PS uspokojivé, začnite hrať skupiny po 2 tremolách, možno s krátkou prestávkou medzi skupinami. Potom zvýšte tempo na skupiny po 3 a potom po 4 tremolách. Najlepší spôsob, ako zrýchliť tremolá, je striedať tremolá a oktávy. Zrýchlite oktávu a skúste prejsť na tremolo pri tejto vyššej rýchlosti. Teraz už len stačí striedať ruky a budovať si vytrvalosť. Opäť platí, že budovanie vytrvalosti nie je ani tak budovanie svalov, ako skôr vedieť, ako sa uvoľniť a ako používať správne pohyby. Oddeľte ruky od tela; nezväzujte systém ruka-paže-telo do jedného tuhého uzla, ale nechajte ruky a prsty fungovať nezávisle od tela. Mali by ste voľne dýchať, bez toho, aby vás ovplyvňovalo to, čo robia prsty. \emph{\textbf{Pomalé cvičenie s prehnanými pohybmi je prekvapivo účinné, takže sa k nemu vráťte vždy, keď narazíte na problém.}}

Pre pravú ruku (oktáva Bb, takt 149) by mala byť 1 hlasnejšia ako 5, ale pre obe ruky by mali byť tichšie tóny jasne počuteľné a ich zrejmým účelom je zdvojnásobiť rýchlosť v porovnaní s hraním oktáv. Nezabudnite cvičiť potichu; neskôr môžete hrať hlasnejšie, kedykoľvek budete chcieť, keď si osvojíte techniku a vytrvalosť. Je dôležité vedieť hrať potichu a zároveň počuť každú notu pri najrýchlejších rýchlostiach. Cvičte, kým pri konečnej rýchlosti nebudete môcť hrať tremolá dlhšie, ako v skladbe potrebujete. Konečný efekt ľavej ruky je konštantný rev, ktorého hlasitosť môžete modulovať nahor a nadol. Nižšia nota poskytuje rytmus a horná nota zdvojnásobuje rýchlosť. Potom cvičte stúpajúce tremolá podľa pokynov v notovom zápise.

Časť Grave, ktorá začína túto prvú vetu, nie je jednoduchá, hoci tempo je pomalé, kvôli svojmu nezvyčajnému rytmu a rýchlym beham v taktoch 4 a 10. Rytmus prvého taktu nie je jednoduchý, pretože chýba prvá nota druhej doby. Aby ste sa naučili správny rytmus, použite metronóm alebo pri nácviku pravej ruky doplňte jednotlivé rytmické noty s ľavou rukou (LH). Hoci je rytmus 4/4, je jednoduchšie, ak zdvojnásobíte noty ľavej noty a cvičíte ich ako 8/8. Beh v takte 4 je veľmi rýchly; v poslednej skupine 1/128 nôt je 9 nôt; preto sa musia hrať ako trioly, dvojnásobným tempom ako predchádzajúcich 10 nôt. To si vyžaduje 32 nôt na dobu, čo je pre väčšinu klaviristov nemožné, takže možno budete musieť použiť rubato; \emph{\textbf{správnym tempom môže byť podľa pôvodného rukopisu polovica uvedeného}}. 10. takt obsahuje toľko nôt, že v Doverovej edícii zaberá 2 riadky! Posledná skupina 16 nôt v tempe 1/128 sa opäť hrá dvojnásobnou rýchlosťou ako predchádzajúcich 13 nôt, čo je pre väčšinu klaviristov neuveriteľne rýchlo. Štvortónový chromatický prstoklad (III.5h) môže byť pri takýchto rýchlostiach užitočný. Každý študent, ktorý sa túto časť Grave učí prvýkrát, musí starostlivo počítať noty a doby, aby získal jasnú predstavu o tom, o čo ide. Tieto šialené rýchlosti môžu byť chybou redaktora.

Prvá (a tretia) veta je variáciou na tému z Grave. Táto slávna téma „Dracula`` bola prevzatá z ľavej dolnej časti prvého taktu; je zrejmé, že ľavá dolná časť nesie emocionálny obsah, hoci pravá dolná časť nesie melódiu. Venujte pozornosť tvrdému staccatu a sf v taktoch 3 a 4. V taktoch 7 a 8 sa posledné tóny troch stúpajúcich chromatických oktáv musia hrať ako

Noty 1/16, 1/8 a 1/4, ktoré v kombinácii so stúpajúcou výškou a crescentom vytvárajú dramatický efekt. Toto je pravý Beethoven s maximálnym kontrastom: ticho-hlasno, pomaly-rýchlo, akordy s jednou notou a zložité akordy. V Beethovenovom rukopise nie je žiadna indikácia pedálu.
\end{document}